\title{跋}

历时二十日，终于写完了。

你现在学会 \CONTEXT\ 了吗？

我还没有学会。

这份文档所介绍的 \CONTEXT\ 知识或许尚不及它全部功能的百分之一，且 \CONTEXT\ 依然在发展着。学无止境，以生之有涯，随知识之无涯，这样不好。弱水三千，仅取一瓢饮，会好一些。

我应该去外面走走了，不然春天又悄悄离去了。

\blank[line]
\rightaligned{2023 年 4 月}

\midaligned{\hl[8]}
\blank[line]
\indentation 用了大约十二天，更新了这份文档。大概有三分之一的内容是重写的，诸如命令行环境、列表、表格、插图、数学公式、幻灯片、参考文献等内容，基本是重新写了一遍，最后一章\boxquote{杂技}是新增的，其余章节在局部上有所增删。

此次更新，原本的计划里是想加入几节 Lua 与 \CONTEXT\ 协作的内容，但是后来又放弃了，原因是这部分内容对于常规的文档排版而言并非必须。不过，我在知乎网站上专门为 \CONTEXT\ 创建了一个专栏，其地址在

\useURL[zhuanlan][https://zhuanlan.zhihu.com/c_1932476336884156084]
\midaligned{\boxquote{\from[zhuanlan]}}

\noindent 该专栏不仅有一些关于 \CONTEXT\ 的文章，也有一些文章试图以实践者的方式讲述 MetaPost 语言。

\blank[line]
\rightaligned{2025 年 9 月}

\midaligned{\hl[8]}
\blank[line]
\useURL[expolorer-cc][https://github.com/Explorer-cc][][@Explorer-cc]
\indentation 这次修订工作，主要是修正了文档里的一些书写错误，并对之前部分颇为草率的字句予以细化，主体内容基本上未有增减。在此需要特别感谢 \from[expolorer-cc] 的悉心匡正，否则我不会有太多动力去完成这些工作。

对文档的版权略作了修改，补充了一些简单的条款，以满足 Debian 社区对 TeXLive 软件包的授权方面的严谨需求，而我本人对这些条款所能起到的作用实际上一无所知……要感谢 AI 比我懂这方面的事，所补充的条款也是它帮我写的。

第 \in[Figure] 章的示例中的插图，我换成了我童年时的偶像——多啦 A 梦。

第 \in[Installation] 章的结语部分有较大修改之处。今年风靡一时的养龙虾事件，想必已经让很多人意识到了，命令行并非落后的玩意，反而颇为有用。此外，第 \in[misc] 章的结语也做了修改，为新增加的第 \in[programming] 章做好了铺垫。

第 \in[programming] 章是一份简单的 \TEX\ 编程入门指南，原本我是计划用它呈现 \TEX\ 编程在排版工作中的一些应用，但是在经过一些尝试后，我觉得还是让它尽量简单一些，仅立意于打消你对 \TEX\ 编程的畏惧。

\blank[line]
\rightaligned{2026 年 4 月}
